% DOCUMENT CLASS %
\documentclass[a4paper, 12pt, oneside, ngerman, twocolumn]{article}   % for A4 Paper, 12pt fontsize.
\setlength{\columnsep}{1cm}

% PACKAGES%
\usepackage{amsfonts}   % For a basic mathfont (many will be replaced by stix)
\usepackage{amsmath}    % For basic math symbols
%\usepackage{bbm}        % For \mathbbm (better version of \mathbb)
\usepackage{mathrsfs}   % For \mathscr
\usepackage{hyperref}   % For footnotes
\usepackage{csquotes}   % For \textquote{}
\usepackage{IEEEtrantools}  % For better alignments
\usepackage{tikz}   % For a macro
\usepackage{listings} % For code
\usepackage{xcolor}   % For ... code
\usepackage[babel]{microtype}
\usepackage{enumitem}
\usepackage[left=1.5cm,right=1.5cm,top=2cm,bottom=2cm]{geometry}
\usepackage[color=gray]{todonotes}
\usepackage{epigraph}   % For epigraphs
\usepackage{fancyhdr} % For headers and footers

% To control how (section) titles look.
\usepackage{titlesec}
\titleformat{\section}
{\normalfont\Large\scshape\bfseries}
{\thesection}
{1em}
{}

% Language
\usepackage[ngerman]{babel}

% If needed
\usepackage[nohyperlinks,smaller]{acronym}

% For tests
\usepackage{lipsum}

%%% FONT CONFIGURATION %%%
\usepackage{fontspec}

% Updated version of Crimson Text
\setmainfont{Cochineal}
\setmonofont{Cascadia Code}[Scale=0.8]
\usepackage[cochineal]{newtxmath} % To use Cochineal as a Math Font
%%%

% Citing
\usepackage[style=alphabetic, backend=biber]{biblatex}
\addbibresource{citations.bib}

% MAKE DESCRIPTION ENVIRONMENTS HAVE ITALICIZED ITEMS %
\renewcommand{\descriptionlabel}[1]{\hspace{\labelsep}\textit{#1}}

% SETS %
\newcommand*{\R}{\mathbb R}    % Set of real numbers
\newcommand*{\N}{\mathbb N}    % Set of natural numbers
\newcommand*{\Z}{\mathbb Z}    % Set of integers
\newcommand*{\Q}{\mathbb Q}    % Set of rational numbers

% BASIC CONFIGURATION %
\pagestyle{plain}
\allowdisplaybreaks

% CODE %
\definecolor{codepurple}{HTML}{7f0055}
\definecolor{comment}{RGB}{0,128,0}
\definecolor{number}{RGB}{9,136,90}
\definecolor{string}{RGB}{163,21,21}

% Things that are the same in all styles
\lstdefinestyle{basic}{
    basicstyle=\ttfamily\footnotesize,
    breakatwhitespace=false,
    breaklines=true,
    frame=leftline,
    keepspaces=true,
    numbers=left,
    firstnumber=1,
    numbersep=5pt,
    showspaces=false,
    showstringspaces=false,
    showtabs=false,
    tabsize=2,
    commentstyle=\itshape\color{comment},
    stringstyle=\color{string},
    keywordstyle=\bfseries\color{codepurple},
    texcl=true,
    mathescape=true
}

\lstdefinestyle{pseudo}{
    style=basic,
    language=python,
    morekeywords={do, od, for, if, fi, then, else, to, def, in, forall, exists, while, accept, reject},
}

\lstset{style=pseudo}
